\documentclass[output=paper,colorlinks,citecolor=brown]{langscibook}
\ChapterDOI{10.5281/zenodo.20541850}
\author{Elena Smirnova\orcid{0000-0003-3763-9951}\affiliation{Université de Neuchâtel}}
\title{German \emph{werden}-future: A case of failed grammaticalization?}
\abstract{In present-day German, the construction \textit{werden} + infinitive is usually considered as an analytic future tense form, and the diachronic development of this construction is described as a typical case of grammaticalization. However, the synchronic status of this construction is still a matter of discussion. On the one hand, the function of this construction is not restricted to expressing futurity, but encompasses modal and evidential functions. On the other hand, when it comes to marking futurity, \textit{werden} + infinitive is not obligatory. Although the relatively high grammaticalization status of the construction seems uncontested, especially when compared with the less grammaticalized copula verb \textit{werden} from which it originated, this paper will argue that the grammaticalization process has not come to full completion and was halted at some point in time. What happened then is that the use of the futurate \textit{werden} + infinitive was gradually restricted to a rather small range of particular contexts of use, i.e. a functional niche, where it unambiguously serves as a future tense marker.}
%move the following commands to the "local..." files of the master project when integrating this chapter


%\bibliography{localbibliography}
%\newcommand{\orcid}[1]{}
% \usepackage{graphicx} %package to manage images
% \graphicspath{ {./images/} }
% \usepackage[justification=centering]{caption}


\IfFileExists{../localcommands.tex}{
   \addbibresource{../localbibliography.bib}
   \usepackage{tabularx,multicol}
\usepackage{url}
\urlstyle{same}

\usepackage[english]{babel}
 
\usepackage{langsci-optional}
\usepackage{langsci-lgr}
\usepackage{langsci-gb4e}

\usepackage{langsci-branding}

%\usepackage{tabularx}

% elena
\usepackage{graphicx} %package to manage images
\usepackage{paralist}

% sarah & lena
\usepackage[justification=centering]{caption}
\usepackage{rotating}

% robin
\usepackage{lscape}
\usepackage{tikz}
\usetikzlibrary{tikzmark}

% dominika
\usepackage{multirow}
\usepackage{hanging}

% olaf & dylan
\usepackage{subcaption}

% svetlana
\usepackage{verbatim}
\usepackage{enumitem}

% wiemer
%\usepackage[markup=underlined]{changes}
%\usepackage{hyphenat}

% wir
\usepackage{longtable}

% for crossref in the introduction jiabin
\usepackage{xr}
   % \newcommand*{\orcid}{}
%

\makeatletter
\let\thetitle\@title
\let\theauthor\@author
\makeatother

\newcommand{\togglepaper}[1][0]{
   \bibliography{../localbibliography}
   \papernote{\scriptsize\normalfont
     \theauthor.
     \titleTemp.
     To appear in:
     E. Di Tor \& Herr Rausgeberin (ed.).
     Booktitle in localcommands.tex.
     Berlin: Language Science Press. [preliminary page numbering]
   }
   \pagenumbering{roman}
   \setcounter{chapter}{#1}
   \addtocounter{chapter}{-1}
}
%
\newbool{bookcompile}
\booltrue{bookcompile}
\newcommand{\bookorchapter}[2]{\ifbool{bookcompile}{#1}{#2}}
%
% %\newcommand{\orcid}[1]{}
\graphicspath{ {./figures/} }
%
\let\eachwordone=\itshape


% Seongha
% \newcommand{\koreex}[1]{{\fontspec{Baekmuk Batang}{#1}}}

% Kutida
\newfontfamily\thaifont{Noto Sans Thai} 
\newcommand{\textthai}[1]{{\thaifont #1}} % für Thai


\newfontfamily\chinesefont{Noto Sans CJK SC}
\newcommand{\textchinese}[1]{{\chinesefont #1}} % Für Chinesisch

% wiemer
\def\hyph{-\penalty0\hskip0pt\relax}
% \definechangesauthor[color=NavyBlue]{SH}



\newcolumntype{K}{>{\raggedright\arraybackslash}p{0.6\textwidth}}  % Right column with 0.7 width


% \newfontfamily\krn[Scale=MatchLowercase,BoldFont=SourceHanSerif-Bold.otf]{SourceHanSerif-Regular.otf}
% \AdditionalFontImprint{Source Han Serif KO}
% \XeTeXlinebreaklocale 'ko'
% \renewcommand{\krn}{}



% Rhee
\newfontfamily\koreanfont{Noto Sans CJK KR} 
\newcommand{\textkorean}[1]{{\koreanfont #1}} % für Koreanisch

\renewcommand{\REF}[2][]{(\ref{#2}#1)}
   %% hyphenation points for line breaks
%% Normally, automatic hyphenation in LaTeX is very good
%% If a word is mis-hyphenated, add it to this file
%%
%% add information to TeX file before \begin{document} with:
%% %% hyphenation points for line breaks
%% Normally, automatic hyphenation in LaTeX is very good
%% If a word is mis-hyphenated, add it to this file
%%
%% add information to TeX file before \begin{document} with:
%% %% hyphenation points for line breaks
%% Normally, automatic hyphenation in LaTeX is very good
%% If a word is mis-hyphenated, add it to this file
%%
%% add information to TeX file before \begin{document} with:
%% \include{localhyphenation}

\hyphenation{
    con-tri-bute
    par-a-digm
    pe-ri-phra-ses
    achieve-ment
    accomp-lish-ment
    We-li-nchen-kang-ci-kok
    cross-ling-u-is-tic
    rec-ord-s
    trun-ca-ted
    de-ri-va-tion-al
    stem-de-ri-va-tion-al
    pre-sent
    in-flu-en-ced
    sys-tem
    ques-tions
    Bor-kovs-kij
    Da-ny-len-ko
    e-qui-va-lents
    re-cog-nized
    In-ter-fe-renz-schicht
    Aus-ei-nan-der-se-tz-ung
    Früh-neu-hoch-deutsch-kor-pus
    Ak-tions-art
    Smir-no-va
    Whi-le
    Ar-chive
    bi-ginnan
    Rei-chen-bach
    ger-ma-nis-ti-sche
    Re-fe-renz-kor-pus
    le-xi-ka-lisch-en
    Alt-hoch-deutsche
    Li-te-ra-tur
    dia-chron
    Dia-lek-to-lo-gie
}

\hyphenation{
    con-tri-bute
    par-a-digm
    pe-ri-phra-ses
    achieve-ment
    accomp-lish-ment
    We-li-nchen-kang-ci-kok
    cross-ling-u-is-tic
    rec-ord-s
    trun-ca-ted
    de-ri-va-tion-al
    stem-de-ri-va-tion-al
    pre-sent
    in-flu-en-ced
    sys-tem
    ques-tions
    Bor-kovs-kij
    Da-ny-len-ko
    e-qui-va-lents
    re-cog-nized
    In-ter-fe-renz-schicht
    Aus-ei-nan-der-se-tz-ung
    Früh-neu-hoch-deutsch-kor-pus
    Ak-tions-art
    Smir-no-va
    Whi-le
    Ar-chive
    bi-ginnan
    Rei-chen-bach
    ger-ma-nis-ti-sche
    Re-fe-renz-kor-pus
    le-xi-ka-lisch-en
    Alt-hoch-deutsche
    Li-te-ra-tur
    dia-chron
    Dia-lek-to-lo-gie
}

\hyphenation{
    con-tri-bute
    par-a-digm
    pe-ri-phra-ses
    achieve-ment
    accomp-lish-ment
    We-li-nchen-kang-ci-kok
    cross-ling-u-is-tic
    rec-ord-s
    trun-ca-ted
    de-ri-va-tion-al
    stem-de-ri-va-tion-al
    pre-sent
    in-flu-en-ced
    sys-tem
    ques-tions
    Bor-kovs-kij
    Da-ny-len-ko
    e-qui-va-lents
    re-cog-nized
    In-ter-fe-renz-schicht
    Aus-ei-nan-der-se-tz-ung
    Früh-neu-hoch-deutsch-kor-pus
    Ak-tions-art
    Smir-no-va
    Whi-le
    Ar-chive
    bi-ginnan
    Rei-chen-bach
    ger-ma-nis-ti-sche
    Re-fe-renz-kor-pus
    le-xi-ka-lisch-en
    Alt-hoch-deutsche
    Li-te-ra-tur
    dia-chron
    Dia-lek-to-lo-gie
}
   \boolfalse{bookcompile}
   \togglepaper[4]%%chapternumber
}{}

\begin{document}
\maketitle


\section{Introduction}\label{sec:smirnova:1}
It is rather uncontroversial that present-day German does not have a default future construction in the sense that (i) there exists a verbal periphrasis with a dedicated future auxiliary and (ii) this periphrasis is highly grammaticalized and (iii) it is obligatory in all contexts with future time reference. Though most grammar books of present-day German present \emph{werden} + infinitive as a future tense construction (see \ref{ex:smirnova:1}) along with other analytic tenses such as perfect or pluperfect, this construction in fact only marginally functions as an analytical future tense. 

\ea \label{ex:smirnova:1}
\gll Er wird im Urlaub sein.\\
     he werden.\textsc{3sg}.\textsc{pres} on vacation be.\textsc{inf}\\
\glt `He \textsc{wird} be on vacation.'\\
    \ea \label{ex:smirnova:1a}
    He will be on vacation. = future time reference 
    \ex \label{ex:smirnova:1b}
        He might be on vacation. = epistemic assessment
    \ex \label{ex:smirnova:1c}
        He is bound to be on vacation. = evidential assessment 
    \z
\z

Without going into much detail on the longstanding debate regarding semantic and functional characteristics of \emph{werden} + infinitive in present-day German, it is fair to say that it is not an easy task to determine the function of this construction. In the literature, there are three main accounts that may be briefly summarized as follows:
\par One branch of researchers, the so-called “temporalists” (see e.g. \cite{MatzelUlvestad1982}; \cite{Thieroff1992}) consider \emph{werden} to be primarily a future tense auxiliary (see the paraphrase in \ref{ex:smirnova:1a}) and derive other readings from this future semantics. 
\par The so-called “modalists” (see esp. \cite{vater1975}) on the other hand, treat \emph{werden} essentially as a modal verb, namely as an epistemic modal verb, see the interpretation given in (\ref{ex:smirnova:1b}); and other readings are then derived from the modal semantics of this verb.
\par A third option is to see \emph{werden} as primarily evidential in meaning (see esp. \cite{DiewaldSmirnova2010}) namely an inferential evidential verb, as in (\ref{ex:smirnova:1c}). 
\par The present paper does not attempt to answer the question of what counts as the central or core meaning of the construction \emph{werden} + infinitive. Instead, it focuses on the temporal reading of this construction and argues that this reading is restricted to a very marginal set of contexts which constitute a functional niche of this construction. The goal of this paper is first, to show that the grammaticalization of \emph{werden} + infinitive has not been proceeding further since late Early New High German (ENHG), and second, to explore the functional niche it occupies in present-day German.

\par The paper is structured as follows. In Section \ref{sec:smirnova:2}, I will briefly recall some known facts about the diachrony of several future tense candidates in the history of German and motivate my research questions. In Section \ref{sec:smirnova:3}, I will report the results of a diachronic corpus study that concentrates on some particular contexts of use, namely subordinate temporal and conditional clauses. It is in these contexts that the hampered grammaticalization of the future marking manifests itself prominently. In Section \ref{sec:smirnova:4}, I will focus on synchronic corpus data in zooming in on another context of use, namely the combination of \emph{werden} and the infinitive verb \emph{sein} `to be', in order to pinpoint the characteristics of the functional niche the construction occupies in present-day German. The paper will conclude with some thoughts on the reasons why \emph{werden} has not reached an advanced degree of grammaticalization in German.

\section{Future constructions in the diachrony of German} \label{sec:smirnova:2}
In present-day German, \emph{werden} + infinitive is used to mark posteriority, that is, to locate the described event in the future relative to the speech time, as for example in (\ref{ex:smirnova:2}). 

\ea \label{ex:smirnova:2}
\gll Bereits am 1. April wird das Strandbad Müggelsee seine Pforten öffnen.\\
     already on 1st April werden.\textsc{3sg}.\textsc{pres} the Lido Müggelsee its gate.\textsc{pl} open.\textsc{inf}\\
\glt `The Müggelsee Lido will open as early as April 1st.'\\
\hfill [Berliner Tageblatt 1931, DWDS-Kernkorpus]
\z

From the diachronic perspective, it is well known that \emph{werden} + infinitive was not the only construction on the way to becoming a future marker in the history of German. Prior and partly in parallel to \emph{werden} + infinitive, especially in Middle High German (MHG), the combination of \emph{werden} and present participle was used to mark future time reference (cf. \ref{ex:smirnova:3}; see \cite{Hartmann2021}; \cite{Ihden2023}).

\ea \label{ex:smirnova:3}
\gll Salic ſit ir die nu weinint. wan ir werdint lachinde. \\ 
     blessed be.\textsc{2.pl.pres} you who now weep.\textsc{2pl}.\textsc{pres}. when you werden.\textsc{2pl}.\textsc{pres} laugh.\textsc{pres.ptcpl.pl}\\
\glt `Blessed are you who weep now, for you shall laugh.'\\
\hfill [REM\footnote{REM stands for \emph{Referenzkorpus Mittelhochdeutsch}, a corpus of Middle High German data (see \url{https://www.linguistics.rub.de/rem/} for more details).}, 12\textsuperscript{th} c., from \citet[367]{Hartmann2021}]
\z

As regards \emph{werden} + present participle, it became obsolete in the course of the ENHG period, as has been clearly shown by Stefan Hartmann in his diachronic corpus study, cf.: 
\begin{quote}
    Während im REM \emph{werden} + Partizip die deutlich frequentere Konstruktion ist und beide Konstruktionen vom 12. zum 13. Jahrhundert einen Anstieg erfahren, sinkt die Gebrauchsfrequenz der \emph{werden} + Partizip-Konstruktion vom 13. zum 14. Jahrhundert deutlich. [...] \emph{werden} + Infinitiv erlebt einen geradezu explosionsartigen Anstieg, während \emph{werden} + Partizip im eher niedrigen Frequenzbereich verharrt. \citep[377]{Hartmann2021}\footnotemark
\end{quote}\footnotetext{Translation: While \textit{werden} + participle is the considerably more frequent construction, and both constructions experience an increase from the 12\textsubscript{th} to the 13\textsuperscript{th} century, the usage frequency of the \textit{werden} + participle construction decreases considerably from the 13\textsubscript{th} to the 14\textsuperscript{th} century. [...] \textit{werden} + Infinitive experiences an enormous increase, while \textit{werden} + participle remains in a rather low frequency area.}

But also the modal verbs such as \emph{sollen} `shall', \emph{wollen} `want, will', and \emph{müssen} must' were all used with future time reference in the history of German (cf. \ref{ex:smirnova:4}; see \cite{DiewaldHabermann2005} for further details). 

\ea \label{ex:smirnova:4}
\gll Ach Gott hilff / Ach Got hoere / Ach Got / wenn wil es ein end werden? \\
     oh God help.\textsc{imp.sg} / oh God hear.\textsc{imp.sg} / oh God / when will.\textsc{3sg.pres} it a end werden.\textsc{inf}\\
\glt `Oh God help / Oh God hear / Oh God / when will it end?'\\
     \hfill [Dietrich, 16th c., \citet[232]{DiewaldHabermann2005}]
\z

As for the modal verbs, they too did not develop future marking as their central semantics, and did not proceed further on the grammaticalization path towards future constructions in German (cf. \cite[232]{DiewaldHabermann2005}): ``Die Entwicklung von werden zum Futurmarker hat ganz offensichtlich den Grammatikalisierungsprozess der Modalverben zu Futuren gebremst''.\footnote{Translation: The development of \textit{werden} to a future marker has quite obviously curbed the grammaticalization process of modal verbs to futures.}

\par \emph{Werden} + infinitive is hence a clear winner in this diachronic competition with modal verbs and with the construction \emph{werden} + present participle, as only this construction successfully developed into a future tense marker. However, the story of \emph{werden} + infinitive appears to be less successful as it might seem at first glance. 
\par For present-day German, it has been repeatedly observed that this construction is generally very infrequent. To refer to future time, German speakers clearly prefer present tense forms (cf. \citealt{Brons-Albert1982}; \citealt{Hilpert2008}). For example, \citet{Brons-Albert1982} found that only 30\% of all futurate contexts contain \emph{werden} + infinitive. The unmarked expression of future time today is not \emph{werden} + infinitive, but the (futurate) present tense.

\par Thus, it seems that no significant semantic and formal changes have occurred since \emph{werden} + infinitive prevailed over the other constructions in the late Middle High German period, and, most importantly, the frequency of use has remained very low or even decreased since then. In the remainder of this paper, I will argue that this construction has not progressed further on the way towards becoming a highly grammaticalized construction. Instead, it has even regressed, and is now limited to some specialized niches.

\par In the following, I will focus on two local contexts of use of this construction, which I will examine further in order to substantiate my claim: 

\begin{enumerate}
  \item From the diachronic perspective, I will look at conditional and temporal subordinate clauses, i.e. the contexts in which \emph{werden} + infinitive can be unambiguously interpreted as conveying future time reference. Starting from the assumption that the appearance of \emph{werden} + infinitive in conditional clauses represents an advanced step in the grammaticalization process towards a future tense marker, a detailed exploration into these contexts will give us valuable insights into the diachronic development of this construction.
  \item From the synchronic perspective, the contexts with the infinitive verb \emph{sein} `be' will be dealt with in a more detailed quantitative and qualitative analysis of corpus data from present-day German. As previously shown in studies by \citet{Hilpert2008} and \citet{Hartmann2021}, and as will be confirmed by my own diachronic analysis in Section \ref{sec:smirnova:3}, this verb is the most frequent infinitive verb that combines with \emph{werden}. Given the fact that \emph{sein} itself is a highly polyfunctional verb in present-day German, the mere co-occurrence of the two does not provide us with any substantial information about the functional and/or semantic range of \emph{werden} + \emph{sein}. The in-depth analysis of these contexts presented in Section \ref{sec:smirnova:4} will identify the relevant characteristics of the semantic niche \emph{werden} + infinitive occupies in present-day German.
\end{enumerate}

\section{\emph{Werden} + infinitive in conditional and temporal clauses} \label{sec:smirnova:3}

In a comprehensive diachronic corpus study of \emph{würde} + infinitive in MHG and ENHG (\cite{Smirnova2006}), I also looked at the indicative form of this construction, i.e. \emph{werden} + infinitive. One of the central insights of this study was that, in its diachronic development, \emph{werden} + infinitive passed through an intermediate stage of semantic change where it clearly was used to mark a textually salient sequential relation. \emph{Werden} + infinitive acquired this semantics in the MHG period via the reinterpretation of the original \textit{aktionsart} ingressive meaning (see \cite[259--267]{Smirnova2006} for further details). 
\par In this respect, the development of \emph{werden} + infinitive towards a future tense construction follows a rather unusual trajectory. Starting as a copula verb with aspectual change-of-state semantics (i.e. ingressive aktionsart), it developed into a marker of textual sequence relation before turning into a temporal marker of posteriority. These semantic shifts proceeded incrementally and gradually along a series of syntactic contexts, and they show a particular directionality: (i) from main clauses preceded by subordinate clauses to (ii) independent clauses and (iii) finally to subordinate clauses.
\par In main clauses preceded by subordinate clauses, i.e. in stage (i) of the development, the proposition introduced by \emph{werden} + infinitive is marked as following from some important pieces of information given in the text, as shown for example in (\ref{ex:smirnova:5}) and (\ref{ex:smirnova:6}). 

\ea\label{ex:smirnova:5}
\gll Denn wenn solch stück der Königin Auskopmt zu den weiberen hin, So \textbf{werden} sie verachten gar All ir menner ganz offenbar.\\
	 then when such piece the queen come.out.\textsc{3.sg} to the women to so werden.\textsc{3.pl} they despise.\textsc{inf} even all their men quite obviously \\
\glt `For if such a piece of the queen's head is sent out to the women, they will despise all of their men quite obviously.'\\
\hfill [ENHGCorpus, \citet[262]{Smirnova2006}]
\z

\ea\label{ex:smirnova:6}
\gll Wer nun Tugend / gute Sitten vnd Kuensten liebhat/ \textbf{wird} wissen was er tun sollte. \\
	 who now virtue / good morals and skills like.\textsc{3.sg}/ werden.\textsc{3.sg} know.\textsc{inf} what he do.\textsc{inf} shall.\textsc{3.sg}\\
\glt `Now he who loves virtue / good morals and arts / will know what he should do.'\\
\hfill [ENHGCorpus, \citet[263]{Smirnova2006}]
\z
In \xref{ex:smirnova:5}, \emph{werden} + infinitive occurs in the apodosis of a conditional sentence and introduces a conditional consequence. In \xref{ex:smirnova:6}, on the other hand, \emph{werden} + infinitive introduces an effect of a condition that appears in the form of a free relative clause. Additionally, in both examples \xref{ex:smirnova:5} and \xref{ex:smirnova:6}, this construction can be interpreted as referring to future time. What is important is that in these contexts, the predicates do not refer to a specific future event, but they refer generically to certain types of events. 
\par In the course of ENHG, the temporal future reading of the construction \emph{werden} + infinitive gets gradually emancipated from this kind of textual support and expands into new contexts of use. This is especially the case for uses of the construction in independent clauses, see \xref{ex:smirnova:7}. In these contexts, the marking of textual sequence relation is no longer the function of this construction. 

\ea \label{ex:smirnova:7}
\gll Von dem andern geschriben ist: Dann \textbf{wierst} du \textbf{nemmen} das oppher der gerechtichait. \\
of the other write.\textsc{past.part} be.\textsc{3sg}: then werden\textsc{2.sg} you take.\textsc{inf} the sacrifice the justice\\
\glt `Of the other it is written: Then you will take the sacrifice of justice.'\\
\hfill [ENHGCorpus DR, \citet[261]{Smirnova2006}]
\z

The next and most advanced step in the semantic development of future tense marking is reached when the construction unambiguously bears future time reference and does not bear any textual relation to the preceding text. This reading is most pronounced in the subordinate temporal and conditional clauses, as shown in (\ref{ex:smirnova:8})--(\ref{ex:smirnova:10}). 

\ea \label{ex:smirnova:8}
\gll Dem Lande wird wiederum wohl seyn \textbf{wenn} Lycogenes die Krone tragen \textbf{wird} \\
	 the country werden.\textsc{3.sg} again happy be.\textsc{inf} when Lycogenes the crown wear.\textsc{inf} werden.\textsc{3.sg} \\
\glt `The country will again be happy when/if Lycogenes will wear the crown.'\\
\hfill [ENHGCorpus 147, \citet[301]{Smirnova2006}]
\z

\ea \label{ex:smirnova:9}
\gll \textbf{Wenn} er aber finden \textbf{wird}, daß Jhr Euer Hertz nicht zu ihm lencken wollet, \textbf{nachdem} Jhr von allem ſogenannten Zwange befreyet ſeyn \textbf{werdet}, ſo wird er aufhoͤren, an Euch {zu dencken}. \\
	 when he but find.\textsc{inf} werden.\textsc{3.sg} that you your heart not to him lead.\textsc{inf} want.\textsc{2.pl} after you of all so-called constraint free.\textsc{past.part} be.\textsc{inf} werden.\textsc{2.pl} so werden.\textsc{3.sg} he stop about you to.think.\textsc{inf} \\
\glt `But if he finds that you do not want to turn your heart to him, after you have been freed from all of the aforementioned constraints, he will stop thinking about you.'\\
\hfill [Richardson, 1748, DTA]
\z

\ea \label{ex:smirnova:10}
\gll ſo will ich dasjenige, was der meinigen fehlt, \textbf{ſobald} es moͤglich ſeyn \textbf{wird}, hinzufuͤgen \\
	 so want.\textsc{1.sg} I that what from mine miss.\textsc{3.sg} {as soon as} it possible be.\textsc{inf} werden.\textsc{3.sg} add.\textsc{inf} \\
\glt `I will add what is missing from mine as soon as it will be possible.'\\
\hfill [Gessner, 1745, DTA]
\z
 
It should be noted that the temporal subordinators \emph{nachdem} `after' and \emph{sobald} `as soon as' are much younger than the temporal-conditional \emph{wenn} `when': \empty{\textit{Nachdem}} is first attested since ca. 1550, and \emph{sobald} appears at around 1600, whereas \emph{wenn} `when' was present since the Old High German period. In the context of this study, only the temporal and realis-conditional uses of the temporal-conditional subordinator \emph{wenn} are relevant, as only those are associated with the development of the construction \emph{werden} + infinitive.\footnote{Whereas irrealis-conditional contexts with \emph{wenn} are clearly marked by subjunctive verbal morphology, realis-conditional contexts with indicative verbal form are often ambiguous between temporal and conditional reading throughout the time period.} These subordinators share a general conceptual schema that is relevant for the history of \emph{werden} + infinitive that is investigated in this study. 
\par Essentially, the general temporal schema of these contexts may be conceptualized as follows. The scene consists of two events: the event \textbf{e1} that takes place at time \textbf{t1} and the event \textbf{e2} at time \textbf{t2}. The event \textbf{e1} is expressed in the subordinate clause with the subordinators \emph{wenn} `when', \emph{nachdem} `after', and \emph{sobald} `as soon as', and the event \textbf{e2} is depicted in the main clause that usually follows the subordinate clause. The time \textbf{t1} precedes the time \textbf{t2}, what is often iconically represented in the syntagmatic ordering of the clauses; both events are usually located in the future relative to the speech time. 
\par Importantly, in contexts like (\ref{ex:smirnova:8})--(\ref{ex:smirnova:10}), \emph{werden} + infinitive can only be read as future tense construction. That is, it cannot be read in one of its present readings mentioned in Section \ref{sec:smirnova:1}, i.e. neither as a modal nor as an evidential verb, as these markers are not compatible with subordinate clauses. 

\largerpage[-1]

\par The expansion of \emph{werden} + infinitive into temporal and conditional subordinate clauses along the trajectory presented above thus represents a very advanced step on a grammaticalization path towards future tense marker. From the usual way grammaticalization processes proceed, it is to be expected that in the further stages of this process, the construction \emph{werden} + infinitive should get firmly established in these contexts and become obligatory.
\par However, from what we know about present-day German, \emph{werden} + infinitive is no longer used in these contexts.\footnote{The fact that \emph{werden} + infinitive occurred in these contexts at all is particularly interesting from the cross-linguistic point of view. As is known, `real' future tense constructions behave differently with respect to this type of subordinate clause. The English \emph{will}-future or the synthetic future in French are banned from conditional clauses, whereas in Slavic languages, future markers appear naturally in these contexts.} The following case study is based on data consisting of conditional and temporal clauses, in order to trace recent changes in the development of the construction \emph{werden} + infinitive and see what happened in the time between ENHG and today.


\subsection{Data} \label{sec:smirnova:3.1}
The present study is based on data from three recent diachronic slices of German: 1600–1649, 1700–1749, and 1950–1999. The data comes from two corpora of German: \href{https://www.deutschestextarchiv.de}{Deutsches Textarchiv (DTA)} for the first half of the 17th and of the 18th centuries, and the \href{https://www.dwds.de/d/korpora/kern}{DWDS-Kernkorpus} for the second half of the 20th century, which in my data represents present-day use. Table \ref{tab:smirnova:1} gives an overview of the total occurrences of temporal and conditional clauses in the data. The number of subordinators \emph{wenn} `when/if', \emph{nachdem} `after', and \emph{sobald} `as soon as', tagged in the corpus as KOUS `subordinating conjunction', is taken as a proxy for the number of the clauses they introduce.

\begin{table}
\caption{Subordinate clauses with \emph{wenn}, \emph{nachdem}, and \emph{sobald}}
\label{tab:smirnova:1}
 \begin{tabularx}{\textwidth}{l C R {2cm} R{1.5cm} R{1.5cm} R{1.5cm}}
  \lsptoprule
    period & corpus  & corpus size & \textit{wenn} clauses  & \textit{nachdem} clauses & \textit{sobald} clauses\\
  \midrule
  1600--1649  &   DTA  &     7 203 918 & 9 462      & 347 & 5\\
  1700--1749  &   DTA &     23 500 609 & 64 483     & 7 767 & 437\\
  1950--1999  &   DWDS &  58 068 993 & 88 782 & 5 548 & 2 147\\
  \lspbottomrule
 \end{tabularx}
\end{table}

Figure \ref{fig:smirnova:1} depicts the normalized frequencies per million words of the same types of subordinate clauses. As can be seen, \emph{wenn} is the most frequent subordinator across the three periods, the other two being far less frequent.

\begin{figure}
\includegraphics[width=10cm]{figures/smirnova/Figure_1.pdf}
\centering
\vspace*{-10mm}
\caption{\emph{wenn}, \emph{nachdem}, and \emph{sobald}, normalized frequency}
\label{fig:smirnova:1}
\end{figure}



\par The qualitative analysis is mainly based on the data presented in Table \ref{tab:smirnova:2}, where only subordinate clauses containing the construction \emph{werden} + infinitive are shown. The data amounts to a total of 740 observations that are distributed over three periods of time and the three subordinators under investigation.

\begin{table}
 \caption{Subordinate clauses with \emph{werden} + infinitive}
\label{tab:smirnova:2}
 \begin{tabularx}{.75\textwidth}{l C R{1.5cm} R{1.5cm} R{1.5cm}}
  \lsptoprule
    period & corpus  & \textit{wenn} clauses  & \textit{nachdem} clauses & \textit{sobald} clauses\\
  \midrule
  1600--1649  &   DTA  &      213      & 2 & 0\\
  1700--1749  &   DTA &     312     & 23 & 4\\
  1950--1999  &   DWDS &  162 & 10 & 14\\
  \midrule 
   & & 687 & 35 & 18\\
  \lspbottomrule
 \end{tabularx}
\end{table}

\subsection{Analysis} \label{sec:smirnova:3.2}
Figure \ref{fig:smirnova:2} shows the relative frequency of the clauses containing the construction \emph{werden} + infinitive in the three types of subordinate clauses. 

\begin{figure}
\includegraphics[width=10cm]{figures/smirnova/Figure_2.pdf}
\centering
\vspace*{-10mm}
\caption{Relative frequency of \emph{werden} + infinitive in subordinate clauses with \emph{wenn}, \emph{nachdem}, and \emph{sobald}, in percent}
\label{fig:smirnova:2}
\end{figure}

\par The numbers in Figure \ref{fig:smirnova:2} reveal that the proportions are extremely low: Most bars hardly reach the one percent mark. That means that subordinate clauses with \emph{werden} + infinitive are extremely rare. In the first period, \emph{wenn}-clauses have the highest proportion of \emph{werden} + infinitive (2.25\%). Interestingly, a decline in the use of \emph{werden} + infinitive can be observed in all types of subordinate clauses; it is most pronounced in \emph{wenn}-clauses (dark bars in Figure \ref{fig:smirnova:2}), but a similar progression can be observed for \emph{nachdem}- and \emph{sobald}-clauses. What we see here is a clear tendency for \emph{werden} + infinitive to slowly disappear from the conditional and/or temporal clauses in the course of time. If we look at the younger subordinators \emph{nachdem} and \emph{sobald}, we can see that \emph{werden} + infinitive initially finds its way into these structures, but fails to gain ground and declines in use in the subsequent periods.

\par From this, it can be concluded that \emph{werden} + infinitive gradually loses its contexts of use. In Section \ref{sec:smirnova:3}, conditional and temporal subordinate clauses were described as representing an advanced stage of grammaticalization of \emph{werden} + infinitive as a future tense marker. These contexts were the last into which the construction expanded, and the further grammaticalization process would normally imply further expansion and consolidation of use in these contexts. What we observe, however, is the opposite -- \emph{werden} + infinitive is gradually withdrawing from these contexts. The construction does not proceed further on the grammaticalization path.

\par Furthermore, a closer look into the data reveals that, in the third period, the subordinate clauses differ considerably from each other as concerns the form of the infinitive (see Table \ref{tab:smirnova:3}). 

\begin{table}
\caption{Form of the infinitive, 1950--1999}
\label{tab:smirnova:3}
 \begin{tabularx}{.85\textwidth}{l R{1.5cm} R{1.5cm} R{1.5cm}}
  \lsptoprule
    subordinator & \textit{wenn}  & \textit{nachdem}  & \textit{sobald}\\
  \midrule
  \emph{werden} + infinitive  &   137  &      4      & 1\\
  \emph{werden} + infinitive perfect &   25 &     10     & 9\\
  \midrule 
   \% infinitive perfect & 15\% & 71\% & 90\%\\
  \lspbottomrule
 \end{tabularx}
\end{table}

Whereas \emph{wenn}-clauses still predominantly contain \emph{werden} + infinitive (see e.g. \ref{ex:smirnova:11}), \emph{nachdem}- and \emph{sobald}-clauses show a clear preference for \emph{werden} + infinitive perfect, usually called future perfect or \emph{Futur II} in German (see e.g. \ref{ex:smirnova:12}).

\ea \label{ex:smirnova:11}
\gll Wenn wir trotzdem an diesen ungesicherten Allgemeinbegriffen \textbf{festhalten} \textbf{werden}, so liegt darin mehr als ein unentschuldbarer methodologischer Leichtsinn.\\
    when we nonetheless to these unsecured  generalities cling.\textsc{inf} werden.\textsc{3.pl} so lie.\textsc{3.sg} there more than a inexcusable methodological recklessness \\
\glt `If we will nevertheless cling to these unsecured generalities, there is more than an inexcusable methodological recklessness in this.'\\
\hfill [Dahrendorf, 1957, DWDS-Kernkorpus]
\z

\ea \label{ex:smirnova:12}
\gll Sie hat versprochen sich zu entscheiden, sobald sie das Leichentuch für ihren Schwiegervater Laertes \textbf{fertiggewebt} \textbf{haben} \textbf{wird} \\
    she has promised \textsc{refl} to decide {as soon as} she the shroud for her father-in-law Laertes finish.\textsc{past.part} have.\textsc{inf} werden.\textsc{3.sg} \\
\glt `She has promised to decide as soon as she will have finished weaving the shroud for her father-in-law Laertes.'\\
\hfill [Schwanitz, 1999, DWDS-Kernkorpus]
\z

The form \emph{werden} + perfect infinitive is usually employed to mark a completion of an event in the future and is particularly tied to contexts of temporal clauses such as those analyzed here. 

\par Another important observation concerns the verb form in the main clauses. The relative proportion of different verb forms is depicted in Figure \ref{fig:smirnova:3}.

\begin{figure}
\includegraphics[width=10cm]{figures/smirnova/Figure_3.pdf}
\centering
\vspace*{-10mm}
\caption{Relative frequency of different verb forms in main clauses, in percent}
\label{fig:smirnova:3}
\end{figure}

\par What we see in Figure \ref{fig:smirnova:3} is a clear tendency towards the decline of the form \emph{werden} + infinitive across the three periods: Whereas in the 17th century, \emph{werden} + infinitive was the dominant form in both clauses of the sentence, in the 20th century, \emph{werden} + infinitive in the main clause amounts only to 40\% of all verb forms. On the other hand, the rise of the use of simple present forms is particularly pronounced: The use of the present tense increases from ca. 10\% in the first period to 40\% in the third period. That is, main clauses follow the same general tendency we have observed above, and the construction \emph{werden} + infinitive gets gradually replaced by the present tense forms, both in main and in subordinate clauses of temporal and conditional sentences.

\par The data presented so far shows that \emph{werden} + infinitive gradually disappears from temporal and conditional clauses with \emph{wenn}, \emph{nachdem}, and \emph{sobald} (see Figure \ref{fig:smirnova:2}). In present-day German, \emph{werden} + infinitive is almost completely absent from the temporal clauses with \emph{nachdem} and \emph{sobald} (see Table \ref{tab:smirnova:3}); \emph{werden} + perfect infinitive (= ``future perfect'') is used there instead. As seen from Figure \ref{fig:smirnova:2}, \emph{werden} appears first in these contexts, but fails to gain ground and disappears gradually from them. Moreover, this downward tendency holds for main clauses of these sentences as well (see Figure \ref{fig:smirnova:3}).

\par A last observation to be drawn from the analysis of this dataset concerns the infinitive verbs combining with \emph{werden}.\footnote{For the infinitive verbs, only the forms with the present infinitive, and not those with the perfect infinitive, were analyzed.} Table \ref{tab:smirnova:4} gives an overview of the ten most frequent infinitive verbs in each period. 

\begin{table}
\caption{Infinitive verbs, top ten collocates}
\label{tab:smirnova:4}
 \footnotesize 
 \begin{tabularx}{\textwidth}{ll R{.35cm} ll R{.35cm} ll R{.35cm}}
  \lsptoprule
    \multicolumn{3}{c}{1600–1649} & \multicolumn{3}{c}{1700–1749}  & \multicolumn{3}{c}{1950–1999}\\
  \midrule
  verb & gloss & N & verb & gloss & N & verb & gloss & N \\
  \midrule
  \emph{kommen} & come  & 33 &   \emph{kommen} & come  & 41 & \emph{sein} & be & 64\\
  \emph{erscheinen} & appear  & 18 &  \emph{sein} & be & 28 & \emph{haben} & have & 9\\
  \emph{sein} & be & 16 & \emph{haben} & have & 19 & \emph{kommen} & come & 4\\
  \emph{anschauen} & look & 5 & \emph{erscheinen} & appear & 13 & \emph{gelingen} & succeed & 3\\
  \emph{offenbaren} & reveal & 5 & \emph{sehen} & see & 9 & \emph{lesen} & read & 3\\
  \emph{sagen} & say & 5 & \emph{werden} & become & 8 & \emph{werden} & become & 3\\
  \emph{sehen} & see & 5 & \emph{hören} & hear & 7 & \emph{ausmachen} & turn out & 2\\
  \emph{werden} & become & 5 & \emph{posaunen} & broadcast & 7 & \emph{führen} & lead & 2\\
  \emph{angehen} & concern & 4 & \emph{aufhören} & stop & 5 & \emph{gehen} & go & 2\\
  \emph{brechen} & break & 4 & \emph{offenbaren} & reveal & 5 & \emph{heiraten} & marry & 2\\
 \lspbottomrule
 \end{tabularx}
\end{table}

\par As can be seen from Table \ref{tab:smirnova:4}, some of the infinitive verbs, such as \emph{kommen} `come' or \emph{haben} `have', have combined with \emph{werden} from the beginning; others have become more or less associated with the construction in the course of time. In the verb list of the last period, however, \emph{sein} `be' is not only the most frequent verb, but it is also almost the only verb in this combination. With 64 occurrences, it stands in sharp contrast to other infinitive verbs, which are attested in the data with very rare occurrences. 
\par \citet[385]{Hartmann2021} analyzed the usage of the construction \emph{werden} + infinitive in the ENHG period and found that towards the end of the period, i.e. in the 17th century, the most frequent infinitive verb in this construction is \emph{sein} (23 occurrences in his data), followed by \emph{haben} `have' (with 8 occurrences). In a study based on data from the 19th and 20th century, \citet[148]{Hilpert2008} detected the same verbs, that is, \emph{sein} `be' (n=172), \emph{kommen} `come' (n=62), \emph{haben} `have' (n=40), and \emph{werden} `become' (n=40) among the most frequent infinitive verbs in this construction. In each study, it is thus the verb \emph{sein} that turns out to be the most frequent infinitive verb that combines with \emph{werden}. This suggests that the changes detected above and that concern the temporal and conditional clauses are not only characteristic of these particular contexts, but manifest a general tendency in the development of the construction \emph{werden} + infinitive. 

\par In sum, the results of the diachronic corpus study show that in the last four centuries, the construction \emph{werden} + infinitive has not progressed further on its way towards becoming a grammaticalized future tense marker: 
\begin{enumerate}
    \item With regard to usage frequency, there is a decrease in use and a gradual replacement of \emph{werden} + infinitive by the futurate present tense. As concerns its use in temporal and conditional subordinate clauses, after initial uses, it failed to expand further, and this tendency holds also for independent uses of this construction. That is, no further context-expansion of the construction in the sense of \citet{Himmelmann2004} can be observed since ENHG. Instead, the contexts of use become more and more restricted. 
    \item With regard to the infinitive verbs used in the construction, the range of regularly used verb types decreased, most of them being used rarely in the construction. Importantly, \emph{werden} + infinitive established a particular preference for the infinitive verb \emph{sein} `be'. This limitation to one dominant usage pattern with \emph{sein} clearly goes against the tendency of obligatorification \citep{Lehmann1995}. During the advanced stages of a grammaticalization process, obligatorification would manifest itself in the increasing token and type frequencies of the infinitive verbs. However, this is clearly not the case for \emph{werden} + infinitive, as shown above.
\end{enumerate} 
To conclude this section, in the light of the evidence presented above, I would like to argue that the development of the construction \emph{werden} + infinitive in the last four centuries constitutes a case of failed grammaticalization. What has happened instead is that \emph{werden} + infinitive has limited itself to some non-frequent and specialized uses. This empirical evidence points to the establishment of a possible functional niche of this construction in modern German. In Section \ref{sec:smirnova:4}, I will deal in more detail with the combination of \emph{werden} + \emph{sein} to see what kind of functional niche this particular pattern occupies in present-day German.

\largerpage[1]

\section{The functional niche of \emph{werden} + infinitive in present-day German} \label{sec:smirnova:4}
Let me start this section by giving some illustrative examples of the combination of \emph{werden} and \emph{sein} (cf. \ref{ex:smirnova:13}--\ref{ex:smirnova:15}).

\ea \label{ex:smirnova:13}
\gll Du \textbf{wirst} mir eines Tages dankbar \textbf{sein}. \\
    you werden.\textsc{2.sg} me one day thankful be.\textsc{inf} \\
\glt `You will be grateful to me one day.'\\
\hfill [Noll, 1996]
\z

\ea\label{ex:smirnova:14}
\gll Was \textbf{wird} inzwischen zu Hause alles geschehen \textbf{sein}? \\
    what werden.\textsc{3.sg} meanwhile at home all happen.\textsc{past.ptcpl} sein.\textsc{inf} \\
\glt `What will have happened at home in the meantime?'\\
\hfill [Grzimek, 1954]
\z

\ea\label{ex:smirnova:15}
\gll Ob und wie dieser Ratschlag funktionierte, darüber \textbf{wird} noch zu sprechen \textbf{sein}.\\
    whether and how that advice work.\textsc{3.sg.past} {about it} werden.\textsc{3.sg} yet to speak.\textsc{inf} be.\textsc{inf} \\
\glt `Whether and how this advice worked, will have to be discussed yet.'\\
\hfill [P. M, 1993]

\z

The examples \xref{ex:smirnova:13}--\xref{ex:smirnova:15} show that the combination of \emph{werden} + \emph{sein} occurs in different syntagmatic configurations and displays an array of possible functions. In \xref{ex:smirnova:13}, \emph{sein} is a copula verb that combines with an adjective (\emph{dankbar} `grateful'); the combination with \emph{werden} renders future time reference. In \xref{ex:smirnova:14}, on the other hand, \emph{sein} is a perfect tense auxiliary modifying the full verb \emph{geschehen} `happen', and the combination \emph{werden} + \emph{past participle} + \emph{inf} constitutes a case of Future~II or future perfect. This combination may refer to future time, but mostly, it is interpreted as a modal construction, i.e. expressing the epistemic uncertainty of the speaker. In \xref{ex:smirnova:15}, \emph{sein} appears as a part of another construction, i.e. \emph{sein} + \emph{zu}~`to'-infinitive with the modal meaning of obligation. The combination with \emph{werden} can be interpreted as referring to an obligation that will hold at a later point in time, or as a textual lookahead, where the author announces that something will (have to) be treated later in the text.
\par Given this variety of functions the syntagmatic combination of \emph{werden} and \emph{sein} may carry in individual contexts, it is worthwhile to inspect the individual contexts in more detail in order to see whether there are patterns in the data that point to the existence of a particular functional niche of the pattern \emph{werden} + \emph{sein}. This is the aim of the following synchronic case study.

\subsection{Data} \label{sec:smirnova:4.1}
The present analysis is based on a random sample of 1000 observations that stem from the most recent slice of the \href{https://www.dwds.de/d/korpora/kern}{DWDS-Kernkorpus20} \citep{geyken2007dwds}, i.e. from the years 1950–1999.\footnote{The search string used to extract the examples of the combinations of \emph{werden} and \emph{sein} from the corpus was the following: /werde|wirst|wird|werden|werdet/ \&\& ``sein with \$p=/V.*INF/'' \#asc\_date[1950-00-00, 1999-99-99]; i.e. I searched for all present indicative forms of the verb \emph{werden} in combination with the infinitive form of \emph{sein} within a sentence. The obtained hits were then randomized,  manually inspected for false positives and annotated according to the variables described in this section.} The data were manually annotated with respect to the type of the element the infinitive verb \emph{sein} combines with. The following categories were used for annotation:
\begin{compactitem}
    \item \textsc{adjective} (see \ref{ex:smirnova:13} above)
    \item \textsc{noun phrase} (see \ref{ex:smirnova:16})
    \item \textsc{prepositional phrase} (see \ref{ex:smirnova:17})
    \item \textsc{adverb} (see \ref{ex:smirnova:18})
    \item \textsc{zu-infinitive} (see \ref{ex:smirnova:15} above) 
    \item \textsc{past participle} (see \ref{ex:smirnova:14} above)
    \item \textsc{(adjectival) past participle} (see \ref{ex:smirnova:19})
\end{compactitem}



\ea\label{ex:smirnova:16}
\gll Das ausführende Organ der Bank wird \textbf{der} \textbf{Vorstand} sein. \\
    the executive body {of the} bank werden.\textsc{3.sg} the {board of directors} be.\textsc{inf} \\
\glt `The executive body of the bank will be the board of directors.'\\
\hfill [AdG 40, 1970]
\z


\ea\label{ex:smirnova:17}
\gll Ich bin zu Hause, wie gesagt, und solange es hell ist, werde ich \textbf{im} \textbf{Garten} sein. \\
    I be.\textsc{1.sg} at home as said and {as long as} it light be.\textsc{3.sg} werden.\textsc{1.sg} I {in the} garden be.\textsc{inf} \\
\glt `I am at home, as I said, and as long as it is light outside, I will be in the garden.'\\
\hfill [Born, 1979]
\z


\ea\label{ex:smirnova:18}
\gll Peter wird natürlich auch \textbf{da} sein. \\
    Peter werden.\textsc{3.sg} {of course} too there be.\textsc{inf} \\
\glt `Peter will of course be there too.'\\
\hfill [Horster, 1950]
\z


\ea \label{ex:smirnova:19}
\gll Dann, und nur dann, wird ein erster Schritt auf dem langen Wege zur nuklearen Abrüstung \textbf{getan} sein. \\
    then and only then werden.\textsc{3.sg} a first step on the long way to nuclear disarmament do.\textsc{past.ptcpl} be.\textsc{inf} \\
\glt `Then, and only then, will a first step on the long road to nuclear disarmament be taken.'\\
\hfill [AdG 28, 1958]
\z

In combinations with adjectives (\ref{ex:smirnova:13}), noun phrases (\ref{ex:smirnova:16}), prepositional phrases (\ref{ex:smirnova:17}), and adverbs (\ref{ex:smirnova:18}), \emph{sein} qualifies as a copula verb. For this reason, these combinations will be treated further below as one sub-group of \emph{sein}-usages. As noted above, the combination of \emph{sein} with \emph{zu} `to'-infinitive qualifies as a construction on its own, and therefore it will be treated as such in the following. In the contexts like \xref{ex:smirnova:14}, \emph{sein} will be considered a perfect auxiliary.

In combinations like (\ref{ex:smirnova:19}), \emph{sein} is usually classified as passive auxiliary, and the construction \emph{sein} + past participle is called stative passive or \emph{Zustandspassiv} (cf. e.g. \cite{Helbig1983}; \cite[567--568]{DUDEN2016}). However, the grammatical classification of the construction as stative passive is not uncontroversial. Without going into much detail at this point, I will follow here the account of \citet{Maienborn2007} who has convincingly shown that the past participle in this construction behaves like an ordinary adjective, and \emph{sein} in these combinations can be considered a ``normal" copula verb. For this reason, I will treat such combinations as pertaining to the sub-category of copula verb \emph{sein} mentioned above. Table \ref{tab:smirnova:5} gives the frequency of these elements in the sample.

\begin{table}%[H]
\caption{Complements of \emph{sein}-infinitive}
\label{tab:smirnova:5}
 \begin{tabularx}{.45\textwidth}{l R{1.5cm}}
  \lsptoprule
    complement & N \\
  \midrule
  adjective  &  388\\
  noun phrase  &   264\\
  \textit{zu}-infinitive & 85\\
  adjectival participle & 81\\
  prepositional phrase & 81\\
  adverb & 54\\
  past participle & 47\\
 \midrule 
   TOTAL & 1000\\
   \lspbottomrule
 \end{tabularx}
\end{table}

Figure \ref{fig:smirnova:4} summarizes the combinations of \emph{werden} and the infinitive of \emph{sein} according to the three categories of the verb \emph{sein} introduced above:
\begin{compactitem}
    \item copula verb in combination with adjectives, adjectival past participles, noun phrases, prepositional phrases, and adverbs
    \item modal construction \emph{sein} + \emph{zu} `to'-infinitive and
    \item perfect auxiliary in combination with past participles of full verbs.
\end{compactitem}

\begin{figure}%[H]
\includegraphics[width=10cm]{figures/smirnova/Figure_4.pdf}
\centering
\vspace*{-10mm}
\caption{Frequency of different \emph{sein}-variants in the data}
\label{fig:smirnova:4}
\end{figure}

As can be already seen from Table \ref{tab:smirnova:5} and Figure \ref{fig:smirnova:4}, the distribution of different variants of the infinitive verb \emph{sein} is extremely skewed.\footnote{Even if I had considered the 81 cases of so-called \emph{Zustandspassiv} as a separate category, this would not have considerably changed the distribution.} In the next section, I will look more closely into the functional range of these combinations. 

\subsection{Analysis} \label{sec:smirnova:4.2}
In the preceding section, three categories of the verb \emph{sein} were distinguished according to the elements the verb combines with, and to the semantics the whole syntagmatic configuration with \emph{werden} conveys. Let me start by discussing the two categories that are only sparsely presented in the data and make up 13.2\% of the analyzed sample. 
\subsubsection{The future perfect form \emph{werden} + past participle + \emph{sein}} \label{sec:smirnova:4.2.1}

\largerpage[1]

\par The future perfect form is attested 47 times in the sample. As mentioned in \tabref{tab:smirnova:5} and \figref{fig:smirnova:4}, in this context \emph{werden} functions as a modal-evidential verb, and it cannot be considered exclusively a future tense auxiliary. In this reading, \emph{werden} behaves similarly to other epistemically used modal verbs like \emph{müssen} `must' or \emph{dürfte} `may' (see \ref{ex:smirnova:20a}--\ref{ex:smirnova:20b} and \ref{ex:smirnova:21a}--\ref{ex:smirnova:21b}):

\ea \label{ex:smirnova:20}
\ea \label{ex:smirnova:20a}
\gll Es \textbf{wird} wohl 'ne größere Taube gewesen sein.\\
    it werden.\textsc{3.sg} probably a larger pigeon be.\textsc{past.ptcpl} be.\textsc{inf} \\
\glt `It will probably have been a larger pigeon.'\\
\hfill [Bieler, 1983]

\ex \label{ex:smirnova:20b}
\gll Es \textbf{muss/dürfte} wohl 'ne größere Taube gewesen sein \\
it must.\textsc{3.sg}/may.\textsc{3.sg.subj} probably a larger pigeon be.\textsc{past.ptcpl} be.\textsc{inf} \\
\glt `It must have been a larger pigeon.'
    \z
\z

\ea \label{ex:smirnova:21}
\ea \label{ex:smirnova:21a}
\gll Für die ``Verräter" \textbf{wird} dieser Rausschmiß nicht unerwartet gekommen sein.\\
    for the traitors werden.\textsc{3.sg} this expulsion not unexpectedly come.\textsc{past.ptcpl} be.\textsc{inf} \\
\glt `For the ``traitors" this expulsion will not have come unexpectedly.'\\
\hfill [ZEIT, 1954]

\ex \label{ex:smirnova:21b}
\gll Für die ``Verräter" \textbf{muss/dürfte} dieser Rausschmiß nicht unerwartet gekommen sein \\
for the traitors werden.\textsc{3.sg} this expulsion not unexpectedly come.\textsc{past.ptcpl} be.\textsc{inf} \\
\glt `For the ``traitors" this expulsion should not have come unexpectedly.'
    \z
\z

These combinations are very low in frequency (see Figure \ref{fig:smirnova:4} in Section \ref{sec:smirnova:4.1}). Moreover, they are functionally similar to other means of expressing epistemic and/or evidential assessment in present-day German. Furthermore, future time reference, i.e. future in the past, constitutes only one possible interpretation of these combinations, whereas the modal nuances are much more pronounced. 
\par This suggests that this usage of \emph{werden} + \emph{sein} does not represent a special functional niche of the more abstract construction \emph{werden} + infinitive. Rather, the future perfect forms attested in the analyzed sample should be seen in a broader context of the usage of perfect infinitive forms in present-day German, including perfect infinitive forms with the perfect auxiliary \emph{haben}, which in many contexts of use display very similar formal and functional characteristics.

\largerpage[-2]

\subsubsection{\emph{Werden} and the modal construction \emph{sein} + \emph{zu} `to'-infinitive} \label{sec:smirnova:4.2.2}
The combination of \emph{werden} with the modal construction \emph{sein} + \emph{zu} `to'-infinitive amounts to 85 observations in the analyzed sample (= 8.5\%). What makes these contexts particularly interesting is the fact that the combination with \emph{werden} unambiguously renders a future tense reading. Compare for example (\ref{ex:smirnova:22a}) and (\ref{ex:smirnova:23a}) with (\ref{ex:smirnova:22b}) and (\ref{ex:smirnova:23b}) respectively, where \emph{werden} has been removed, and only the present tense form of \emph{sein} + \emph{zu} `to'-infinitive is present.

\ea \label{ex:smirnova:22}
\ea\label{ex:smirnova:22a}
\gll Es \textbf{wird} nie wieder gutzumachen sein. \\
    it werden.\textsc{3.sg} never ever {to.make.good} be.\textsc{inf} \\
\glt `It will never be possible to make up for it.'\\
\hfill [Bachmann, 1971]

\ex\label{ex:smirnova:22b}
\gll Es \textbf{ist} nie wieder gutzumachen.\\
it be.\textsc{3.sg} never ever {to.make.good}.\textsc{inf} \\
\glt `It can never be made up for.'
    \z
\z

\ea\label{ex:smirnova:23}
\ea\label{ex:smirnova:23a}
\gll Ihre Aufenthaltserlaubnis \textbf{wird} jedes Jahr zu erneuern sein, und ihre Arbeitserlaubnis \textbf{wird} jedes Mal zu erneuern sein, wenn ... \\
    their {residence permit} werden.\textsc{3.sg} every year to renew.\textsc{inf}  be.\textsc{inf}, and their {work permit} werden.\textsc{3.sg} every time to renew.\textsc{inf}  be.\textsc{inf} when ...\\
\glt `Their residence permit will have to be renewed every year, and their work permit will have to be renewed every time when ...'\\
\hfill [AdG 39, 1969]

\ex\label{ex:smirnova:23b}
\gll Ihre Aufenthaltserlaubnis \textbf{ist} jedes Jahr zu erneuern, und ihre Arbeitserlaubnis \textbf{ist} jedes Mal zu erneuern, wenn ...\\
their {residence permit} be.\textsc{3.sg} every year to renew.\textsc{inf}, and their {work permit} be.\textsc{3.sg} every time to renew.\textsc{inf} when ...\\
\glt `Their residence permit must be renewed every year, and their work permit must be renewed every time when ...'
    \z
\z

In (\ref{ex:smirnova:22b}) and (\ref{ex:smirnova:23b}), the future tense reading gets lost, and only the modal semantics of the construction \emph{sein} + \emph{zu} `to'-infinitive remains, which in turn may be read as a general, tenseless predication. In opposition to (\ref{ex:smirnova:22a}) and (\ref{ex:smirnova:23a}), however, it is clear that the added value of \emph{werden} consists of conveying future tense semantics. That is, the modal semantics of the construction \emph{sein} + \emph{zu} `to'-infinitive is compositionally combined with the future tense semantics of \emph{werden}. Thus, in these contexts, we find \emph{werden} in clear opposition to the simple present tense, and in these contexts, the simple present is not able to produce clear futurate readings.
\par To sum up so far, the combination of \emph{werden} and a conventionalized construction \emph{sein} + \emph{zu} `to'-infinitive with its own modal passive meaning serves to unambiguously mark the future time reference of the proposition. In these contexts, \emph{werden} + infinitive is not functionally equivalent to the present tense, as the present tense form of the construction does not allow for a future tense interpretation. I will return to this point in Section \ref{sec:smirnova:4.3}.

\subsubsection{\emph{Werden} and the copula verb \emph{sein}} \label{sec:smirnova:4.2.3}
The great majority of observations of \emph{werden} + \emph{sein} contains the infinitive of the copula verb \emph{sein}, i.e. 868 observations (= 86.8\%). Only in a small fraction of these observations, namely in 12 sentences, does \emph{werden} not primarily evoke a future tense interpretation, but may be attributed a modal-evidential reading (see \ref{ex:smirnova:24}–\ref{ex:smirnova:25}). Note that in these contexts, modal adverbials such as \emph{wohl} `surely/probably' are present in the text and contribute as well to the modal meaning of the proposition.

\ea\label{ex:smirnova:24}
\gll Was bei Mikrowellen, Joysticks und Zahnbürsten erlaubt ist, \textbf{wird} doch wohl bei der Soziologie nicht verboten \textbf{sein}. \\
    what for microwaves joysticks and toothbrushes allow.\textsc{past.ptcpl} be.\textsc{3.sg} werden.\textsc{3.sg} \textsc{mp}{\footnotemark} surely for the sociology not forbidden be.\textsc{inf}\\
    \footnotetext{German modal particles (\textsc{mp}) are notoriously difficult to translate. For this reason, I resort to the label \textsc{mp} in the glosses.}
\glt `What is allowed for microwaves, joysticks and toothbrushes will surely not be forbidden for sociology.'\\
\hfill [ZEIT, 1996]
\z

\ea\label{ex:smirnova:25}
\gll Das \textbf{wird} doch nicht etwa schon die Polizei \textbf{sein}? \\
    that werden.\textsc{3.sg} MP not MP already the police be.\textsc{inf} \\
\glt `Surely that won't be the police already?'\\
\hfill [Grass, 1959]
\z

In the vast majority of cases, however, the future temporal interpretation is clearly the dominant reading present in the context. Sometimes (in 196 out of 856 or in ca. 23\%), temporal adverbials or other explicit contextual information with future time interpretation are present, as in (\ref{ex:smirnova:17})–(\ref{ex:smirnova:19}) and (\ref{ex:smirnova:26}), but this certainly does not need to be the case (see \ref{ex:smirnova:16} and \ref{ex:smirnova:27}). 


\ea\label{ex:smirnova:26}
\gll Im Jahr 2005 \textbf{wird} Sprache im Umgang mit dem Computer in jeder Hinsicht verfügbar \textbf{sein}. \\
    in year 2005 werden.\textsc{3.sg} language in handling with the computer in every aspect available be.\textsc{inf} \\
\glt `In 2005, language will be available in every aspect of computer use.'\\
\hfill [BILD, 1998]
\z

\ea\label{ex:smirnova:27}
\gll Viele potentielle Produzenten \textbf{werden} nicht {in der Lage} \textbf{sein}, das erforderliche Kapital zu erhalten. \\
    many potential producers werden.\textsc{3.pl} not able be.\textsc{inf} the necessary capital to obtain.\textsc{inf}\\
\glt `Many potential producers will not be able to obtain the necessary capital.'\\
\hfill [AdG 46, 1976]
\z

It may be concluded that for the combinations of \textit{werden} with the copula verb \textit{sein}, the function of \textit{werden} unambiguously and almost exclusively consists of marking future time reference. It is in these contexts that the future meaning of the construction \textit{werden} + infinitive is highly distinctive from the other contexts of use. 
\par The frequency-based evidence presented so far is already very telling: (i) \textit{Sein} is by far the most frequent infinitive verb that combines with \textit{werden} in present-day German (see Section \ref{sec:smirnova:3.2}); (ii) it is the copula verb \textit{sein} that mostly combines with \textit{werden} (see Section \ref{sec:smirnova:4.1}). This suggests that the combination of \textit{werden} and the copula verb \textit{sein} constitutes a very significant use of \textit{werden} + infinitive and may be qualified as a niche. The qualitative evidence presented in this section certifies that in these contexts, \textit{werden} functions as a future tense marker. In other words, the future tense reading of \textit{werden} + infinitive is largely restricted to the combinations of \textit{werden} and \textit{sein}, whereas in all other contexts, it is the (futurate) present tense that is employed to mark future tense reference. In Section \ref{sec:smirnova:4.3}, I will present a tentative account of why combinations with the copula \textit{sein} represent a suitable context for future uses of \textit{werden} + infinitive in present-day German. 

\subsection{Discussion} \label{sec:smirnova:4.3}
It is generally aknowledged that German has three copula verbs: \textit{sein}, \textit{werden}, and \textit{bleiben}. In present-day German, the copula verb \emph{sein} is a stative copula used to ascribe qualities to all kinds of referents, and it stands in semantic opposition to two further copula verbs, i.e. to the inchoative \emph{werden} on the one hand, and to the durative stative \emph{bleiben} on the other, as in (\ref{ex:smirnova:28}).

\ea\label{ex:smirnova:28}
\gll Lena ist/wird/bleibt böse\textsubscript{Adj}/meine Freundin\textsubscript{NP}.\\
    Lena be./werden./remain.\textsc{3.sg} angry/my friend \\
\glt `Lena is/becomes/remains angry/my friend.'
\z

In the following, only the pair \emph{sein} vs. \emph{werden} will be relevant for our purposes. The semantic difference between the two is one of aktionsart: Whereas \emph{sein} is stative and atelic, \emph{werden} is inchoative and telic. These copula verbs are usually seen to behave very similarly with respect to the predicative elements they combine with -- that is, they may combine with adjectives, with noun phrases, with prepositional phrases, and with adverbs (see Section \ref{sec:smirnova:4.1}). 
\par However, there is an important difference in syntactic behaviour between the two. The range of the usage contexts of the copula \emph{werden} is more restrictive than those of \textit{sein}. Whereas \emph{sein} may be combined with all types of predicative elements, irrespective of their semantics, the use of \textit{werden} seems to be restricted to complements that express incremental changes. That is, the use of \textit{werden} is ruled out in combination with locative expressions such as \textit{im Garten} `in the garden' or with non-gradable adjectives such as \textit{tot} `dead'. Table \ref{tab:smirnova:6} summarizes these differences. 

\begin{table}
 \caption{Complements of \emph{sein}- vs. \emph{werden}-copula}
\label{tab:smirnova:6}
\small
 \begin{tabularx}{\textwidth}{llll}
  \lsptoprule
   complement & \emph{sein} & \emph{werden} & gloss\\
  \midrule
  gradable property  &  A ist böse. & A wird böse. & `angry'\\
  non-gradable property & A ist tot. & *A wird tot. & `dead'\\
  location & A ist im Garten. & *A wird im Garten. & `in the garden'\\
  location & A ist dort. & *A wird dort. & `there'\\
   \lspbottomrule
 \end{tabularx}
\end{table}

These semantic differences manifest themselves in an interesting way when it comes to future time reference. Combinations with the copula \textit{sein} unambiguously render stative interpretations, and usually locate the ascribed state or property in the present time. This applies to all kinds of expressions, irrespective of their gradable/non-gradable nature, as in (\ref{ex:smirnova:29})--(\ref{ex:smirnova:30}).

\ea \label{ex:smirnova:29}
\gll Er \textbf{ist} (immer) böse, wenn er das hört.\\
he be.\textsc{3.sg} (always) angry when he it hear.\textsc{3.sg} \\
\glt `He is always angry when he hears that.'
\z
\ea\label{ex:smirnova:30}
\gll Sie \textbf{ist} Astronautin.\\
    she be.\textsc{3.sg} astronaut \\
\glt `She is an astronaut.'
\z

Combinations with the copula \textit{werden} with complements that refer to gradable properties or imply incremental changes, on the other hand, are essentially ambiguous between the inchoative reading with present time reference (= ongoing change of state) and the future reading (= result state in the future). It is only by means of further contextual elements such as temporal adverbials that a non-ambiguous present or future reading is achieved, as in (\ref{ex:smirnova:31a}--\ref{ex:smirnova:31b}).

\ea\label{ex:smirnova:31}
\ea\label{ex:smirnova:31a}
\gll Er \textbf{wird} langsam böse. \\
    he werden.\textsc{3.sg} slowly angry \\
\glt `He is getting angry.'

\ex\label{ex:smirnova:31b}
\gll Er \textbf{wird} böse, wenn er das hört. \\
    he werden.\textsc{3.sg} angry when he it hear.\textsc{3.sg} \\
\glt `He will be/get angry when he hears that.'
\z
\z

Combinations of the auxiliary \textit{werden} with the copula \textit{sein}, on the other hand, clearly result in a non-ambiguous future tense reading, as in (\ref{ex:smirnova:32}):

\ea\label{ex:smirnova:32}
\gll Er \textbf{wird} böse sein. \\
    he werden.\textsc{3.sg} angry be.\textsc{inf}\\
\glt `He will be angry.'
\z

For the expressions that refer to non-gradable properties or non-incremental changes, the situation is rather different. They may be combined with the copula \textit{sein} to refer to a present state,\footnote{The copula verb \emph{sein} in its present tense form is in principle suited for futurate contexts and can be used to refer to future states, especially if some adverbial modification is present in the text that unambiguously refers to future, cf. \emph{Morgen ab 16 Uhr \textbf{bin} ich {im Garten}} `Tomorrow from 4 p.m. I will be in the garden.' However, temporal adverbials are obligatory in such contexts.} as in (\ref{ex:smirnova:33}). When combined with the auxiliary \textit{werden} and the infinitive of the copula \textit{sein}, they unambiguously refer to a future state, as in (\ref{ex:smirnova:34}). Note that the option with the copula verb \textit{werden} alone is not available (see Table \ref{tab:smirnova:6}). 

\ea\label{ex:smirnova:33}
\ea\label{ex:smirnova:33a}
\gll Sie \textbf{ist} bereit.\\
    she be.\textsc{3.sg} ready \\
\glt `She is ready.'
\ex\label{ex:smirnova:33b}
\gll Sie \textbf{wird} bereit sein.\\
    she werden.\textsc{3.sg} ready be.\textsc{inf}\\
\glt `She will be ready.'
\z
\z

\ea\label{ex:smirnova:34}
\ea\label{ex:smirnova:34a}
\gll Ich \textbf{bin} im Garten.\\
    I be.\textsc{1.sg} in.the garden \\
\glt `I am in the garden.'

\ex\label{ex:smirnova:34b}
\gll Ich \textbf{werde} im Garten \textbf{sein}. \\
I werden.\textsc{1.sg} in.the garden be.\textsc{inf}\\
\glt `I will be in the garden.'
    \z
\z

From what has been observed in previous sections, and from the semantic analysis just presented, it is reasonable to assume that for present-day German copula verbs, there exists a firmly established paradigmatic pair \textit{sein} with present time reference and \textit{werden} + \textit{sein} with future time reference. Though \emph{werden} is essentially an inchoative copula verb and can be used as such in contexts with gradable expressions, its use in the construction \textit{werden} + \textit{sein} has been conventionalized as strong expression of future tense, in opposition to the present tense of the copula \emph{sein}.
\par This paradigmatic relation is strongly pronounced in case of more or less fixed multi-word expressions containing the copula verb \textit{sein}. These include, for example, conventionalized and idiomatic expressions such as \textit{an der Reihe sein} `be next, lit: be in the turn', \textit{in der Lage sein} `be able to, lit: be in the position', \textit{zu Ende sein} `be finished', and the like. The only way to mark future time reference is to employ the future auxiliary \textit{werden}, as in (\ref{ex:smirnova:35})--(\ref{ex:smirnova:36}).

\ea\label{ex:smirnova:35}
\gll Unsere Kultur \textbf{wird} {in der Lage} \textbf{sein}, sich unabhängig von der menschlichen Biologie und ihren Gesetzen zu entwickeln.\\
    our culture werden.\textsc{3.sg} {able} be.\textsc{inf} \textsc{refl} independently of the human biology and its laws to develop.\textsc{inf}\\
\glt `Our culture will be able to develop independently of human biology and its laws.'\\
\hfill [Rechenberg, Informatik, 1991]
\z

\ea\label{ex:smirnova:36}
\gll Und auch hier \textbf{wird} es von entscheidender Bedeutung \textbf{sein}, wie sich die DDR-Führung verhält.\\
    and too here werden.\textsc{3.sg} it of critical significance be.\textsc{inf} how \textsc{refl} the {GDR leadership} behave.\textsc{3.sg}\\
\glt `And here, too, it will be crucial how the GDR leadership behaves.'\\
\hfill [DEDok., 1990]
\z

To sum up, I would like to argue that the future construction \textit{werden} + infinitive have ``survived" in this particular niche -- that is, in combination with the infinitive of copula verb \textit{sein}. In combination with other infinitive verbs, \textit{werden} + infinitive did not develop further, and the futurate present tense was more successful in expressing future time reference. 
\par In this diachronic situation, the contexts with the copula verb \textit{sein} offered a welcome niche for the future construction \textit{werden} + infinitive. And this is, as I believe, at least partly due to the asymmetries of usage contexts of the copula verb \textit{werden} described above. As discussed above, \textit{werden} as a copula cannot be used in present tense contexts with a significant number of complements due to their non-incremental and non-gradable semantics. The future auxiliary in combination with the copula \textit{sein} is, however, fine. It could be hypothesized that the combination of \textit{werden} and \textit{sein} serves as a sort of compensating strategy for this gap. Once established in this particular context, it could survive there and be preserved till the present-day German. Of course, more research is needed to substantiate this idea.

\section{Conclusions} \label{sec:smirnova:5}
In this paper, I have argued that the development of the future construction \textit{werden} + infinitive constitutes a case of failed grammaticalization. Although the relatively high grammaticalization status of the construction seems uncontested in present-day German, especially when semantic and formal characteristics of this construction are concerned, a closer look at the frequency of use and the distribution over different infinitive verbs offers a more differentiated picture. Namely, these factors suggest that the grammaticalization process has not come to its full completion, leaving \textit{werden} + infinitive ``just" an alternative means of expressing futurity. 
\par Not being successful and leaving the grammaticalization path does not need to mean, however, that the construction dies out and disappears. As this study has shown, \textit{werden} + infinitive has survived in a particular context of use, i.e. in combination with the infinitive of the copula verb \textit{sein}. In these contexts, the auxiliary \textit{werden} has become nearly obligatory, and it serves here to mark future time reference. The synchronic motivations of this situation as well as the diachronic details of this process are complex and concern not only the construction \textit{werden} + infinitive itself but also the related copula verb \textit{werden}. In Section \ref{sec:smirnova:4.3}, I offered a first tentative account of this functional niche of \textit{werden} in present-day German. But of course, this idea needs to be followed up in further research and substantiated with more data, synchronically and diachronically.

\il{German} 

\sloppy
\printbibliography[heading=subbibliography,notkeyword=this]
\end{document}
